\section{Configuration des périphériques internes}

\subsection{Timer et génération PWM}

\subsubsection{Timer pour la mesure de distance (TIM4)}
Nous avons configuré le TIM4 pour générer la pulsation de 10 µs nécessaire au déclenchement du capteur HC-SR04 toutes les 250 ms. Comme détaillé dans notre calcul initial:

\begin{lstlisting}[caption={Configuration du TIM4}, label={lst:tim4config}]
htim4.Instance = TIM4;
htim4.Init.Prescaler = 319;
htim4.Init.CounterMode = TIM_COUNTERMODE_UP;
htim4.Init.Period = 65535;
htim4.Init.ClockDivision = TIM_CLOCKDIVISION_DIV1;
htim4.Init.AutoReloadPreload = TIM_AUTORELOAD_PRELOAD_DISABLE;
\end{lstlisting}

Un signal PWM avec un rapport cyclique très court (10µs) est généré sur le canal 2 pour déclencher le capteur ultrasonique.

\subsubsection{Timer pour le servomoteur (TIM3)}
Le TIM3 est configuré pour générer un signal PWM adapté au contrôle du servomoteur:

\begin{lstlisting}[caption={Configuration du TIM3}, label={lst:tim3config}]
htim3.Instance = TIM3;
htim3.Init.Prescaler = 28;
htim3.Init.CounterMode = TIM_COUNTERMODE_UP;
htim3.Init.Period = 59999;
htim3.Init.ClockDivision = TIM_CLOCKDIVISION_DIV1;
htim3.Init.AutoReloadPreload = TIM_AUTORELOAD_PRELOAD_DISABLE;
\end{lstlisting}

Cette configuration génère un signal PWM avec une période de 20 ms (50 Hz), ce qui est la fréquence standard pour contrôler les servomoteurs.

\subsection{Configuration des interruptions EXTI}
Pour mesurer précisément le temps de vol des ondes ultrasoniques, nous avons utilisé les interruptions externes (EXTI) pour détecter les fronts montants et descendants du signal ECHO retourné par le capteur HC-SR04.

\begin{lstlisting}[caption={Callback d'interruption EXTI}, label={lst:exticallback}]
void gpio_trigger_Callback() {
    static uint32_t start_time = 0;
    // Lecture de l'etat du pin
    int8_t echo = HAL_GPIO_ReadPin(SonicSensor_Echo_GPIO_Port, SonicSensor_Echo_Pin);
    if (echo){
        // Enregistrement du temps de debut
        start_time = HAL_GetTick();
    } else {
        // Calcul de la duree
        uint32_t duration = HAL_GetTick() - start_time;
        // Calcul de la distance en cm
        distance = (duration * 340)/100/2;
    }
}
\end{lstlisting}

\subsection{Configuration UART}
La communication série est configurée pour permettre l'envoi de la distance mesurée et la réception des commandes de changement de mode ou de position du servomoteur.

\begin{lstlisting}[caption={Initialisation UART}, label={lst:uartinit}]
MX_USART2_UART_Init();
HAL_UART_Transmit(&huart2, (uint8_t *)"Hello World\r\n", 13, HAL_MAX_DELAY);
\end{lstlisting}

